\documentclass[a4paper,12pt,oneside]{report}
\usepackage[T1]{fontenc}
\usepackage[utf8]{inputenc}
\usepackage[ngerman]{babel}
\usepackage{graphicx}
\usepackage{changepage}
\usepackage[headheight=21.75pt]{geometry}
\usepackage{hyperref}
\usepackage[printonlyused]{acronym}
\usepackage{scrlayer-scrpage}
\usepackage{tocloft}
\usepackage{csquotes}
\usepackage{xcolor}
\usepackage{titlesec}
\usepackage{times}
\usepackage{eulervm}
\usepackage{listings}
\usepackage{subcaption}
\usepackage{float}
\usepackage{amsmath}
\usepackage[patch=none]{microtype}
\usepackage[backend=biber, style=alphabetic, doi=true, url=true]{biblatex}
\addbibresource{bibliography/references.bib}

\usepackage{setspace} % Für den Zeilenabstand
\usepackage{parskip}  % Fügt Abstand zwischen Absätzen hinzu

% set url style
\urlstyle{same}

% Fonts and Styling
\geometry{a4paper, left=35mm, right=30mm, top=25mm, bottom=25mm}
\hypersetup{
    colorlinks=true,
    linkcolor=myblue,     % Farbe der internen Links (z.B. Abkürzungen)
    urlcolor=myred,     % Farbe der URL-Links
    citecolor=mygreen    % Farbe der Zitationslinks
}

% Custom Commands and Colors
% Define the variables used in the title page
\newcommand{\myUni}{Hochschule Darmstadt}
\newcommand{\myFaculty}{Fachbereich Informatik}
\newcommand{\myTitle}{Vergleichende Analyse moderner Waldbrand-Frühwarnsysteme:\\ Evaluation von Synergien zwischen terrestrischer IoT-Sensorik, optischer Überwachung und Satellitendaten}
\newcommand{\myDegree}{Bachelor of Science (B.Sc.)}
\newcommand{\myName}{Maira Schäfer}
\newcommand{\myId}{1132721}
\newcommand{\myProf}{Dr. Johannes Werner}
\newcommand{\myTime}{20. März 2026}

% Custom colors
\definecolor{myred}{HTML}{990000}
\definecolor{myblue}{HTML}{0073c0}
\definecolor{mygreen}{HTML}{008000}
\definecolor{chaptergrey}{RGB}{150,150,150}
\renewcommand{\thepart}{\Roman{part}}
\definecolor{jsonBackground}{rgb}{0.95,0.95,0.95}
\definecolor{jsonKeyword}{rgb}{0.5,0.0,0.0}
\definecolor{jsonString}{rgb}{0.0,0.5,0.0}
\definecolor{jsonComment}{rgb}{0.5,0.5,0.5}

% Configure listings for JSON
\lstdefinestyle{jsonstyle}{
    backgroundcolor=\color{jsonBackground},
    basicstyle=\ttfamily\small,
    keywordstyle=\color{jsonKeyword},
    stringstyle=\color{jsonString},
    commentstyle=\color{jsonComment},
    showstringspaces=false,
    breaklines=true,
    captionpos=b,
    frame=single,
    morestring=[b],
    morecomment=[l][\color{jsonComment}]{\#}
}

% Configure listings for bash
\definecolor{codebg}{rgb}{0.95, 0.95, 0.95}
\lstset{
    backgroundcolor=\color{codebg},
    basicstyle=\ttfamily,
    frame=single,
    breaklines=true,
    language=bash,
    columns=fullflexible
}

% Custom styling
\input{layout/chapter_formatting.tex} % Chapter and title formatting
% Configure header and footer
\clearpairofpagestyles % Löscht alle Voreinstellungen für Kopf- und Fußzeilen
\ohead{\headmark} % Setzt den Kapitelnamen oben rechts
% \cfoot{} % Unten nichts anzeigen

\automark{chapter} % Automatisch den aktuellen Kapitelnamen übernehmen
\ofoot{\pagemark} % Setzt die Seitenzahl unten rechts (falls gewünscht)
\pagestyle{scrheadings} % Verwende scrheadings-Stil für Kopfzeilen
 % Header and footer configuration

\begin{document}
\begin{sloppypar}

\input{layout/title_page.tex}
\input{layout/title_back_page.tex}
\input{chapters/doiw/doiw.tex}
\chapter*{Abstract}
\thispagestyle{empty} % Suppresses the page number

BLABLABLA

\newpage

\chapter*{Zusammenfassung}
\thispagestyle{empty} % Suppresses the page number

BLABLABLA

\clearpage % Ensure the next chapter starts on a new page
\input{layout/table_of_contents.tex}
% Tabellenverzeichnis und Abbildungsverzeichnis anpassen

% Anpassungen für Abbildungsverzeichnis
\chapter*{Abbildungsverzeichnis}
\thispagestyle{empty}
\pagestyle{empty} % Entfernt Seitenzahlen auf den Verzeichnisseiten

% Redefiniere \listoffigures, um den vertikalen Raum zu entfernen
\makeatletter
\renewcommand{\listoffigures}{%
  \begingroup % Änderungen lokal begrenzen
  \parskip=0pt % Kein zusätzlicher Abstand zwischen Absätzen
  \parindent=0pt % Kein Einzug
  \@starttoc{lof} % Start des Abbildungsverzeichnisses ohne zusätzlichen Titelraum
  \endgroup
}
\makeatother

% Manuelles Hinzufügen des Abbildungsverzeichnisses ohne Farbänderungen
\vspace{-\baselineskip} % Reduziert den vertikalen Abstand
\listoffigures

\newpage % Neue Seite beginnen

% Anpassungen für Tabellenverzeichnis
%\chapter*{Tabellenverzeichnis}
%\thispagestyle{empty}
%\pagestyle{empty} % Entfernt Seitenzahlen auf den Verzeichnisseiten

% Redefiniere \listoftables, um den vertikalen Raum zu entfernen
\makeatletter
\renewcommand{\listoftables}{%
  \begingroup % Änderungen lokal begrenzen
  \parskip=0pt % Kein zusätzlicher Abstand zwischen Absätzen
  \parindent=0pt % Kein Einzug
  \@starttoc{lot} % Start des Tabellenverzeichnisses ohne zusätzlichen Titelraum
  \endgroup
}
\makeatother

% Manuelles Hinzufügen des Tabellenverzeichnisses ohne Farbänderungen
\vspace{-\baselineskip} % Reduziert den vertikalen Abstand
\listoftables

\newpage % Neue Seite beginnen

% Seitenstil für nachfolgende Seiten wiederherstellen
\pagestyle{scrheadings}

% Abkürzungsverzeichnis
\chapter*{Abkürzungsverzeichnis}
\thispagestyle{empty}
\pagestyle{empty}

\begin{acronym}[Monochrom]
  \acro{AVHRR}{Advanced Very High Resolution Radiometer}
  \acro{BIRD}{Bispectral Infrared Detection}
  \acro{BIROS}{Bi-spectral Infrared Optical System}
  \acro{IoT}{Internet of Things / Internet der Dinge}
  \acro{JPSS}{Joint Polar Satellite System}
  \acro{LEO}{Low Earth Orbit / Erdnahe Umlaufbahn}
  \acro{KIWB}{Künstliche Intelligenz-basiertes Warnsystem zur Waldbranderkennung}
  \acro{LoRaWAN}{Long Range Wide Area Network / Langstrecken-Weitbereichsnetz}
  \acro{LPWAN}{Low Power Wide Area Network / Niedrigleistungs-Weitbereichsnetz}
  \acro{MODIS}{Moderate Resolution Imaging Spectroradiometer}
  \acro{Monochrom}{Schwarz-weiß Sensor}
  \acro{NIR}{Near-Infrared / Nahinfrarot}
  \acro{Suomi NPP}{Suomi National Polar-orbiting Partnership}
  \acro{SWIR}{Short-Wave Infrared / Kurzwelliges Infrarot}
  \acro{TET-1}{Technologieerprobungsträger 1}
  \acro{UAV}[UAV]{Unmanned Aerial Vehicle}
  \acro{VIIRS}{Visible Infrared Imaging Radiometer Suite}
  \acro{VOCs}{Volatile Organic Compounds / Flüchtige organische Verbindungen}
 
  
  
  
  
  
 


\end{acronym}
\pagestyle{scrheadings}

\setcounter{page}{0} % Reset the page counter. Start counting here
\part{Einleitung}\label{pt:einleitung}
\addtocontents{toc}{\protect\setcounter{tocdepth}{-1}} % Schaltet das TOC kurzzeitig "stumm"
\chapter{Einleitung}\label{ch:einleitung}             % Erzeugt "1 Einleitung" im Text
\addtocontents{toc}{\protect\setcounter{tocdepth}{2}}

Die globale Zunahme von Waldbrandereignissen hat sich in den letzten Jahren von einem saisonalen Phänomen zu einer chronischen, weltweiten Bedrohung entwickelt. Die ökologischen, sozialen und ökonomischen Auswirkungen von Waldbränden sind verheerend: Neben der Zerstörung sensibler Ökosysteme und dem Verlust von Biodiversität führen Waldbrände zur massiven Freisetzung von Treibhausgasen, was die anthropogenen Erderwärmung weiter verstärkt.
Die Tatsache, dass im Jahr 2024 erstmals die \qty{1,5}{\degree}C-Grenze überschritten wurde, verdeutlicht die Dringlichkeit, denn die zunehmende Austrocknung von Böden und Vegetation begünstigt die Entstehung von Feuern erheblich~\cite{Copernicus2025,Umweltbundesamt2024}.
Längst beschränken sich die wirtschaftlichen Folgen nicht mehr nur auf Regionen wie Kalifornien, wo die staatlichen Budgets bereits im Jahr 2020 um 3,1 Milliarden USD übertroffen wurden~\cite{LegislativeAnalystsOffice2020}. Parallel dazu sind auch in Europa die jährlichen Ausgaben für die Brandbekämpfung auf inzwischen rund zwei Milliarden Euro angestiegen~\cite{JointResearchCentre2023}.

Die entscheidende Variable zur Eindämmung dieser Schäden ist die Zeit. Die Phasen zwischen dem Brandausbruch und der offenen Flamme, die sogenannte Schwelbrandphase, bietet das größte Potenzial für eine effektive Brandbekämpfung mit minimalem Ressourceneinsatz~\cite{TobisJakobs}. Herkömmliche Detektionsmethoden stoßen hier an ihre Grenzen. Während Meldungen durch die Bevölkerung von der lokalen Siedlungsdichte abhängig sind, sind bemannte Feuerwachtürme aufgrund von mangelnder Präzision und arbeitsrechtlicher Bestimmungen kaum noch zukunftsfähig~\cite{DirkSchneider2017}.

Die technische Antwort auf dieses Defizit ist die Implementierung moderner Früherkennungssysteme, die jedoch spezifische Schwachstellen aufweisen. Weltraumgestützte Systeme leiden unter niedrigen Wiederholraten, was eine Echtzeit-Warnung erschwert. Optisch-terrestrische Kamerasysteme sind schneller, reagieren jedoch auf topografische Hindernisse oder Wetterphänomene, was häufig zu ressourcenintensiven Fehlalarmen führt~\cite{Chagger2018}. Internet of Things (\acs{IoT})-basierte Sensoren wiederum versprechen eine präzise Detektion, erfordern jedoch für eine flächendeckende Abdeckung eine kostspielige und wartungsintensive Infrastruktur~\cite{Klick2021}.

Hieraus ergibt sich eine kritische Diskrepanz zwischen den theoretischen Potenzialen einzelner Technologien und den praktischen Anforderungen an ein zuverlässiges System, das eine geringe Latenzzeit mit einer niedrigen Fehlalarmrate verbindet. Ziel der vorliegenden Arbeit ist es daher, diese technologische Lücke durch eine vergleichende Evaluierung zu analysieren. Im Fokus steht die Frage, welche Kombination aus sensorischen und optischen Frühwarnsystemen unter Berücksichtigung von Latenzzeit, Fehlalarmrate und Kosteneffizienz den effektivsten Schutz bietet.

Es wird die Hypothese aufgestellt, dass multimodale Systeme, die lokale chemische Sensoren mit optischen Systemen kombinieren und dabei weltraumgestützte Frühwarnsysteme nutzen, eine signifikant höhere Zuverlässigkeit und Geschwindigkeit bieten als isolierte Einzellösungen.

\part{Theoretische Grundlagen der Branddetektion}\label{pt:grundlagen}
\chapter{Grundlagen der Brandentstehung}\label{ch:grundlagenBrandentstehung}

Die Entstehung eines Brandes setzt das Vorhandensein von drei Komponenten voraus, die im richtigen Mengenverhältnis vorliegen müsssen: brennbares Material, Sauerstoff als Oxidationsmittel und eine Zündtemperatur, die erreicht werden muss, um die chemische Kettenreaktion in Gang zu setzen. Brennbares Material kann dabei in verschiednene Formen auftreten, wie z.B. Holz, Laub oder trockenen Pflanzenreste. Sauerstoff wirkt als Oxidationsmittel und liegt in unserer Atmosphäre in ausreichender Menge vor. Die Zündtemperatur variiert je nach Material und hängt von Faktoren wie Feuchtigkeitsgehalt und Dichte ab. Wenn diese drei Komponenten zusammentreffen, beginnt der Verbrennungsprozess, der als chemische Reaktion der Oxidation eines Brandstoffes durch ein Oxidationsmittel in Gegenwart einer Energiequelle definiert ist. Energiequellen können in diesem Zusammenhang beispielsweise mechanische Reibung, Elektrizität, chemische Reaktionen oder Solarenergie sein.

Ein Waldbrand durchläuft in der Regel mehrere Phasen, beginnend mit der Enstehungsphase, auch als Schwelbrandphase bezeichnet, in der das Feuer noch klein und auf einen begrenzten Raum beschränkt ist. In dieser Phase sind die Materialien oft nur am Glimmen oder es bilden sich erste kleine Flammen, die jedoch noch keine offene Flamme darstellen. Die Hitze- und Rauchentwicklung sind hier noch relativ gering, da die Temperatur zu Beginn noch erträglich ist, allerdings ist der Rauch hochgiftig und sollte in keinem Fall unterschätzt werden. In der Schwelbrandphase lässt sich ein Feuer oft noch mit Wasser, Feuerlöschern oder durch Ersticken löschen. Bleit ein solcher Entstehungsbrand jedoch unbehandelt, kann er sich sehr schnell entwickeln. Sobald das Feuer Nahrung findet, beispielsweise durch andere trockene Materialien, steigt die Temperatur rasant an, der Rauch wird dichter und es kommt zu einem Kronenfeuer, das sich schnell zu einem Vollbrand über mehrere Bäume entwickeln kann, der sich unkontrolliert ausbreitet. Dürreperioden, trockenen Böden und extreme Abholzung verstärken die Brandgefahr und die Ausbreitung des Feuers.


Ökosysteme, in denen Waldbrände zur natürlichen Entwicklung gehören, sind oft an diese Ereignisse angepasst. Dort spielen Brände eine wichtige Rolle bei der Verjüngung von Wäldern, der Freisetzung von Nährstoffen und der Erhaltung der Biodiversität. Im Westen der USA sorgen die Flammen beispielswese dafür, dass die Zapfen sich aufgrund der Hitze öffnen und ihre Samen verteilen, was zur natürlichen Regeneration der wälder betragt. Die Bäume dieser Wälder sind gut durch ihre Rinde geschützt und können so auch größere Brände überstehen. Allerdings geraten auch hier gefördert durch den Klimawandel immer mehr Brände außer Kontrolle, brennen länger, heftiger und größere Flächen als früher.
Besonders gefährdet sind hingegen Wälder, die durch menschliche Eingriffe entstanden sind, wie z.B. Monokulturen oder Plantagen, da diese oft eine geringere Biodiversität aufweisen und damit weniger widerstandsfähig gegenüber Bränden sind. In Deutschland sind beispielsweise Kiefernwälder aufgrund ihrer ätherischen Öle besonders gefährdert, da diese schneller als naturnahe Wälder austrocknen und leicht brennbares Material bilden.


- Natürlcihe Waldbrände durch Blitzeinschalg, aber ehrs selten

- Waldbrände abhängig von Witterung, Brennmatierial und Tageszeit
- Gefährlich, wenn Waldbrände zu heftig, zu lange, zu großflächig an den falschen Orten und zu ungewöhnlcihen Jahreszeiten

- vor Austrieb neuen Grüns steigt Waldbrandgefahr im Frühjahr an und erreicht Ende April/Anfang Mai erstne höhepunkt, Brandgefahr steigt nochmal in Sommermonaten zwischen Ende Juni bis Ende august
- abhängig von Entstehung, Bodenvegetation, Witterung und waldbeständ knnen Brände in vers. Form auftreten
- Bodenfeur am häufigsten, vollfeuer (auch Kronenfeuer genannt) gefährlichste Arten des Brandes -> benötigen zu entstehn und unterhalt immer energiereiches Bodenfeuer
- Bodenfeuer: Brände, die im Laub- und Streuschicht des Waldbodens brennen und relativ langsam voranschreitne
- Kronen- oder wipfelfeuer hingegen verbreiten sich schnell durch Baumkronen und sind schwer zu kontrollieren
- In De nimmt Zahl der Tage mit hohem Brandrisiko zu
- Ursache für Waldbrände zu 96\% von Menschen asgelöst
Lediglich 4,2\% der Brände werden durch natürlcihe Ursache nwie Blitzeinschlag ausgelöst -> erlsöchen oft schnell wieder, da Gewitter mit Blitz in unseren Breiten meist mit starken Neiderschlägen einhergehen
- 20,6\% zwischen 1992 und 2023 lie0 sich Brandstiftung nachweisen, höherre Anteil 23,8\% durch fahrlässtiges Verhalten
- natürliche Waldökosysteme sind robust, Monukulturen anfällig
- Ökologische Folgelast von waldbränden:
    - Entstehung Vielzahl schädlicher emission von Treibhausgase, Feisntaub und anderen Luftschadstoffen bis hin zu Giften wie Kohlenmonoxid, Furanen und Dioxinen
    - Tier- und Pflanzenort vor Ort wird vernichtet
    - Freigebrannter Boden anfälliger für Erosionen, reicht bereits lkeichtet regen
    - Weggespülte Bodenmaterial, wertvoller, humushalter Oberboden fehlt dem Waldökosysstem und ist potzenill gefährlich für nahe Gewässer
    Erodierter Boden büßt Wasserhaltefähigkeit ein ,was Wiederbewaldung erschwert
    Brandfläche sind wasserundurchlässiger und damit trocjenenr als nicht abgebrannte waldbrandflächen
- Dichte Nadelholzwälder unter 40 Jahre als, insbesondere Kifernwälder sind besonders betroffen -> hier besonders trocken udn durch ätherische Öle der Nadelbäume brennen diese beonsder gut
- Beginnt als Bodenfeuer, gesießst drch Reisig, Äste und trockene Pflanzenreste am Waldboden -> Bodenfeuerr begint Nahrung "hinterher" zu laufen
- Wipfelfeuer entsteht -> Springen Flammen nun von Krone zu Kronen spricht man von Vollfeuer oder totalbrand -> wird idR immer durch Bodenfeuer geschützt -> verliert Bodenfeuer seine Kräfte, fällt Vollfeuer zu Boden und sicht sich dort Nahrung

\chapter{Physikalische Detektionsmerkmale}\label{ch:physikalischedetektionsmerkmale}


Brennbare Materialen werden dabei in Brandklassen eingeteilt, die unterschiedliche Brandarten charakterisieren. Wald- oder Buschbrände gehören dabei zur Brandklasse A. Voraussetzung für die Entzündugn von Holz ist eine thermische Aufbereitung, bei der in der ersten Phase Wasserdampf entweicht. Bei weiterem Temperaturanstieg läuft ein Schwelvorgang ab, bei dem zahlreiche Substantzen entstehen, wie z.B. Kohlenstoffmonoxid, Kohlenstoffdioxid, Methan und Wasserstoff. Diese Substanzen brennen utner gelblicher flammenildung und geben dabei so viel Wärme ab, dass die übriggebliebebne Holzkohle in Brand gesetzt wird. Wenn die gasförmigen Stoffe verbrannt sind, gibt es keine Flammebildung mehr und es bleibt nur noch der Glutbrand der Holzkohle übrig.

\chapter{Anforderungen an Frühwarnsysteme/Definition der Vergleichskriterien}\label{ch:anforderungen}

Echtzeitfähigkeit, Fehlalarmresistenz und Robustheit im Gelände

Detektionslatenz, Fehlalarmrate, Skalierbarkeit und ökonomische Effizienz






\part{Aktueller Stand der Technik}\label{pt:standdertechnik}
% chktex-file 18
\chapter{Weltraumgestützte Systeme}\label{ch:weltraumgestuetztesysteme}

Weltraumgestützte Systeme nutzen unterschiedliche Sensorplattformen, Orbit"=Konfigurationen und Instrumente, die je nach Missionsziel spezifische Aufgaben im Rahmen des Waldbrandmanagements übernehmen. Grundsätzlich operieren die meisten dieser Satelliten in erdnahen, oftmals sonnensynchronen Umlaufbahnen (Low Earth Orbit, LEO), um kontinuierlich Landoberflächen streifenweise abzutasten~\cite{TerraLexikon}.

Für die globale Erkennung von aktiven Waldbränden kommen zunehmend moderne Systeme wie die Visible Infrared Imaging Radiometer Suite (VIIRS) zum Einsatz. VIIRS nutzt Sensoren, die auf verschiedenen Satellitenplattformen installiert sind, darunter die Suomi National Polar-orbiting Partnership (Suomi NPP) und das Joint Polar Satellite System (JPSS). Es kombiniert die Aufgaben älterer Sensorgenerationen wie MODIS und AVHRR, die beispielsweise auf dem Terra-Satelliten oder auf der NOAA-POES- und MetOp-Serie installiert waren und bietet eine verbesserte räumliche Auflösung sowie höhere Detektionsgenauigkeit~\cite{MODISLexikon,TerraLexikon}.
VIIRS arbeitet als scannendes Radiometer, welches die Erdoberfläche im sichtbaren und infraroten Spektrum abtastet. Dadurch ist eine großflächige, weltweite und automatisierte Erfassung von Landoberflächentemperaturen und feuerbedingten Anomalien möglich~\cite{VIIRSLexikon}.

Zur Erfassung von Ereignissen mit hohen Temperaturen, wie sie unter anderem bei Waldbränden auftreten, werden Kleinsatelliten-Missionen genutzt. Diese sind spezifisch für den Zweck der Feuerfernerkundung konzipiert worden. Eine Pionierrolle wird hierbei von der deutsche BIRD-Satellit eingenommen, welche weltweit der erste seiner Art war. Auf diesen Erkenntnissen aufbauend wurde die FireBIRD-Mission mit den beiden Satelliten TET-1 (Technologie-Erprobungsträger) und BIROS (Bi-spectral Infrared Optical System) entwickelt. Sie ermöglichen nicht nur die reine Detektion, sondern können auch weitere wichtige Parameter, wie die Feuerintensität und die Charakterisierung der Feuerfront, mit höherer räumlicher Auflösung als vergleichbare Systeme liefern~\cite{DZLR2024}.
Ein zukunftsweisendes Konzept der Kleinsatelliten-Missionen bietet unter anderem das Unternehmen OroraTech, welches acht Satelliten in den Orbit gebracht hat. Die Anordnung dieser Satelliten bietet den entscheidenden Rahmen für zwei wesentliche Innovationen: Die höhere Auflösung des Systems ermöglicht die frühzeitige Erkennung von Bränden bereits ab einer Größe von 4\,$\times$\,4\,m. Dies gewährleistet eine zeitnahe Alarmierung der Einsatzkräfte nach Brandentstehung. Des Weiteren liefern die Satelliten Daten für die Nachmittags- sowie Nachtstunden. Diese Zeiträume wurden bisher von den meisten Systemen kaum abgedeckt.
 Somit lassen sich nun auch Phasen besonders intensiver Brandaktivität lückenlos überwachen. Zur Analyse der enormen Datenmengen und zur Minimierung der Latenzzeit werden die Daten mithilfe künstlicher Intelligenz ausgewertet. Durch das sogenannte Edge-Computing erfolgt die Datenauswertung meist direkt auf dem Satelliten vor der Datenübertragung, sodass die Ergebnisse in Echtzeit vorliegen können~\cite{Seton2025, DZLR2024}.


\chapter{Optisch-terrestrische Systeme}\label{ch:optischTerrestrischeSysteme}

Um den Herausforderungen der Waldbrandfrüherkennung und -überwachung gerecht zu werden, werden neben weltraumgestützten Systemen auch optisch-terrestrische Systeme eingesetzt. Diese bodengebundenen ("terrestrischen") Technologien basieren auf der Verwendung von visueller ("optischer") Kamerasensorik und sind stationär auf ehemaligen Feuerwachtürmen, Gebäuden oder anderen erhöhten Standorten installiert.
Moderne Systeme integrieren eine Sensoreinheit zur Bilddatenerfassung, eine Recheneinheit zur Datenverarbeitung sowie eine KI-basierte Analysesoftware. Neuronale Netze und Deep-Learning-Algorithmen ermöglichen die automatische Erkennung von Rauch- und Feuersignaturen, sodass zielgerichtet und zeitnah gehandelt werden kann~\cite{IQGmbH2023,Winter2021}.

Die Positionierung der Systeme erfolgt mit einer Mindesthöhe von fünf Metern oberhalb der Baumkronen, um ein ungehindertes Sichtfeld zu gewährleisten und somit die frühzeitige Erkennung von aufsteigenden Rauchsäulen zu ermöglichen.
Um eine lückenlose Erfassung zu gewährleisten, sind die Systeme in der Regel mit einem rotierenden Schwenk-Neige-Kopf ausgestattet. Dieser ermöglicht eine Überwachung des Gebietes in einem \qty{360}{\degree}-Radius. Das vom französischen Unternehmen Paratronic entwickelte Erkennungssystem ADELIE (Alert Detection Localisation of Forest Fire) benötigt beispielsweise für eine vollständige Umdrehung maximal zwei Minuten~\cite{Gemmingen2022}.

Hinsichtlich der Sensorik verwenden moderne Systeme hochauflösende, multispektrale Kameras, um unter unterschiedlichen Tages- und Witterungsverhältnissen eine zuverlässige Erkennung zu gewährleisten.
Das ursprünglich vom DLR (Deutsches Zentrum für Luft- und Raumfahrt) für Weltraummissionen entwickelte System "IQ FireWatch" integriert bis zu vier Sensoren in einem Gehäuse. Dazu gehören ein Monochrom-Sensor für die Aufnahme von Bildern am Tag, ein Nah-Infrarot-Sensor (NIR) für Aufnahmen bei Nacht sowie ein RGB-Sensor für zusätzliche Spektralinformationen. Optional kann, vor allem für Industriegelände, ein zusätzlicher thermischer Infrarotsensor integriert werden, der primär auf kleineren Entfernungen eine gesicherte Erkennung von Hitzequellen ermöglicht. Unter optimalen Bedingungen ist das System so in der Lage, Brände bis zu einer Entfernung von 60~km zu erkennen~\cite{IQGmbH2023, Winter2021}.

Die eigentliche Brandfrüherkennung und Automatisierung des Systems erfolgt über eine komplexe Software, die Rohdaten in Echtzeit auswertet. Das Erkennungssystem ADELIE registriert die Bilddaten und vergleicht sie auf Veränderungen, etwa die Verschiebung von Pixeln, die auf Rauch oder Feuer hindeuten könnten. Anschließend werden diese Differenzen unter Zunahme von verschiedenen Algorithmen analysiert, um alle Elemente, die keinen Rauch darstellen, bestmöglich zu eliminieren~\cite{Winter2021}.

Die Analyse moderner Systeme setzt dabei zunehmend auf einen hybriden Ansatz: Klassische, merkmalbasierte Verfahren werden durch Künstliche Intelligenz (KI) ergänzt. Die neuronalen Netzwerke der KI verarbeiten die mehrdimensionalen Bildmatrizen und erlernen komplexe Muster, um tatsächlichen Brandrauch eindeutig von anderen natürlichen Störfaktoren oder Bildveränderungen zu unterscheiden~\cite{IQGmbH2023,Winter2021}.

Wird auf diesem Weg potenzieller Rauch erkannt, schätzt das System die Entfernung zum potenziellen Brandherd ab. Zudem wird geprüft, ob sich der Rauch in einem Bereich befindet, in dem bekannte Rauchquellen, wie Industrieanlagen und Schornsteine, existieren. Dadurch können Fehlalarme minimiert werden. Das System alarmiert die Einsatzkräfte danach jedoch nicht automatisch, sondern übermittelt die Informationen zunächst an die zuständige Leitstelle, damit diese durch einen menschlichen Operator validiert werden können. Die Mitarbeiter nehmen daraufhin selbst mithilfe einer hochauflösenden Kamera mit starkem optischem Zoom eine Prüfung vor, um festzustellen, ob es sich faktisch um einen Brand handelt, bevor sie die Einsatzkräfte alarmieren.
Das Gesamtsystem agiert folglich als intelligentes Entscheidungsunterstützungssystem~\cite{IQGmbH2023, Gemmingen2022}.


\chapter{IoT-basierte Sensornetzwerke}\label{ch:IoTbasierteSensornetzwerke}

Während satelliten- und optisch-terrestrische Systeme oft erst nach der Entstehung einer offenen Flamme erkennen, setzten IoT-basierte Sensornetzwerke direkt auf der Ebene des Waldbodens und damit deutlich früher an. Brände können bereits in der Schwelbrandphase erkannt werden, bevor sich sichtbare Flammen oder Rauchwolken entwickeln. Dafür greift die Technik auf \enquote{Digitale Nasen} (Electronic Noses oder e-Noses) zurück, die versuchen, das biologische Geruchssystem von Menschen und Tieren nachzuahmen, um spezifische Gerüche und luftgetragene chemische Verbindungen zu erkennen und zu identifizieren.
Zu den grundlegenden Komponenten einer digitalen Nase gehören in der Regel Arrays aus chemischen Sensoren, Signalverarbeitungseinheiten und Mustererkennungsalgorithmen. Nimmt eine digitale Nase Brandgase wie Kohlenstoffmonoxid, Wasserstoff oder flüchtige organische Verbindungen (VOCs) wahr, reagieren Sensoren auf diese Stoffe und wandeln die chemischen Informationen in elektrische Signale um. Diese Signale bilden ein komplexes, mehrdimensionales Geruchsprofil, welches von Algorithmen des maschinellen Lernens und künstlichen Intelligenz mit bekannten Profilen von Feuer und Rauch abgeglichen wird~\cite{Chan2024, FayosJordan2024, Kumar2023}.

Ein wegweisendes Beispiel ist das laufende Forschungsprojekt KIWB (KI-basiertes Warnsystem zur Waldbranderkennung). Es setzt bereits in der Entstehung an und hat zum Ziel, Schwelbrände bereits ab einer Ausdehnung von nur einem Quadratmeter sicher zu detektieren. Das System basiert auf dem Einsatz von energieautarken Gassensoren, die über weite Distanzen funkfähig und für den Langzeiteinsatz im Wald ausgelegt sind. Mittels spezialisierter KI-Lösungen sollen die Sensoren spezifische Gassignaturen von Bränden erkennen und von alltäglichen Emissionen unterscheiden. So soll eine zuverlässige Früherkennung ermöglicht werden, die gleichzeitig die Senkung auf eine Fehlalarmquote von unter ein Prozent anstrebt~\cite{BBF2025}.

Im kommerziellen Sektor setzt das Unternehmen Dryad Networks diese Technologie mit seinem Silvanet-System großflächig ein. Das System besteht aus solarbetriebenen, KI-gestützten Gassensoren, die direkt an den Bäumen des Waldes installiert werden und Brände innerhalb weniger Minuten nach der Entzündung erkennen können. Zur Übertragung der Alarmsignale aus unwegsamen Waldgebieten an die Leistelle, wird eine spezielle Funkinfrastruktur genutzt, da klassische Mobilfunknetzte dort oft nicht oder nur schlecht zur Verfügung stehen~\cite{DNG}.

Hierfür kommt das sogenannte Low-Power Wide-Area Network (LPWAN) zum Einsatz. Diese Technologie ist darauf ausgerichtet kleine Datenmengen extrem energieeffizient über große Distanzen zu senden. Der am weitesten verbreitete und lizenzfreie Standard für LPWAN ist LoRaWAN (Long Range Wide Area Network). LoRaWAN ist ein weltweit offenes Kommunikationsprotokoll speziell für energieeffiziente IoT-Geräte und findet Einsatz in den verschiedensten Anwendungsbereichen, von der Landwirtschaft bis hin zur städtischen Infrastruktur. Es ermöglicht die Übertragung von Daten über Entfernungen von mehreren Kilometern, selbst in dicht bewaldeten oder bergigen Gebieten.
Innerhalb dieses Netzwerks übernehmen sogenannte Gateways eine zentrale Rolle. Anstatt direkt mit einem weit entfernten Server zu kommunizieren, senden die Sensorknoten ihre Pakete zunächst an ein nahegelegenes Gateway, welches die empfangenen Daten dann an einen mit fortschrittlicher Technologie ausgestatteten Server weiterleitet. Klassische LoRaWAN-Systeme sind dabei meist sternförmig aufbaut, es gibt jedoch auch Ausnahmen. Das System Silvanet von Dryad Networks nutzt beispielsweise eine Mesh-Topologie, bei der die Daten von den Gateways untereinander weitergereicht werden~\cite{Nodes.sh2024}. 

durch die Kombination aus direkter Gasmessung am Boden, lokaler Datenverarbeitung durch KI und energieeffizienter Funkübertragung bilden IoT-Sensornetzwerke ein wichtiges Element in der Waldbrandfrüherkennung. Sie schließen die Überwachungslücke im kritischen Anfangsstadium eines Brandes und ermöglichen damit eine noch schnellere Reaktion~\cite{Chan2024,DNG}.


\part{Evaluierung und Konzeptentwicklung}\label{pt:evaluierungKonzeptentwicklung}
\chapter{Vergleichende Analyse der Einzelsysteme}\label{ch:vergleichendeAnalyse}

Eine isolierte Betrachtung der etablierten Frühwarnsysteme zeigt, dass jede Technologie über spezifische Stärken verfügt, jedoch auch physikalische oder konzeptionelle Beschränkungen aufweist. Dabei ist zu erkennen, dass alle drei Systeme Unterschiede in Bezug auf die Detektionslatenz, die Fehlalarmresistenz und die Robustheit sowie die Skalierbarkeit und die ökonomische Effizienz aufweisen. Daraus wird ersichtlich, dass keine der Technologien alle Anforderungen optimal erfüllt.

Weltraumgestützte Systeme charakterisieren sich durch ihre globale und großflächige Abdeckung. Satelliten, wie beispielsweise \ac{VIIRS}, verfügen über eine hohe Skalierbarkeit und gestatten die Datenerfassung ohne die Notwendigkeit zusätzlicher physischer Infrastruktur am Boden. Sie liefern  Daten mit gleichbleibender Qualität und fungieren als Grundlage für die Verarbeitung von Informationen aus der ganzen Welt. Gemessen an der überwachten Fläche resultiert dies in einer signifikant hohen ökonomischen Effizienz. Eine Analyse der historischen Entwicklung zeigt, dass in diesem Bereiche größten Schwächen zu verzeichnen waren. Brände wurden von den Satelliten im \acs{LEO} meist erst erkannt, wenn sie bereits eine offene Flamme entwickelt hatten und deshalb eine große thermische Signatur aufwiesen. Kommerziell genutzte Satelliten minimieren diese Latenz mittlerweile durch bordeigenes Edge Computing auf wenige Minuten. Dennoch sind die Systeme durch physikalische Barrieren, wie zum Beispiel dichte Wolkendecken, in ihrer Fehlalarmresistenz und damit in ihrer Zuverlässigkeit eingeschränkt~\cite{DZLR2024, Seton2025,VIIRSLexikon}.

Optisch-terrestrische Systeme, wie IQ FireWatch, bieten demgegenüber eine deutlich höhere Detektionslatenz, da sie ihre Umwelt kontinuierlich in Echtzeit überwachen. Der Einsatz einer multispektralen Sensorik in Kombination mit künstlicher Intelligenz führt zu einer erhöhten Fehlalarmresistenz. Dies ist auf die Fähigkeit zurückzuführen, bekannte Störquellen zu erkennen und zu eliminieren. Für eine lückenlose und landesweite Abdeckung aller Waldflächen, müssten jedoch hunderte Türme mit Strom- und Datenanbindung errichtet werden, was die Kosten und den Wartungsaufwand erheblich steigert. Diese Systeme sind dementsprechend weniger skalierbar und ökonomisch weniger effizient~\cite{Gemmingen2022,Winter2021}.

\acs{IoT}-basierte Sensornetzwerke bieten die geringste Detektionslatenz. Systeme wie \enquote{Silvanet} von Dryad Networks sind in der Lage, Brände bereits in der Schwelbrandphase ab einer Größe von einem Quadratmeter zu erkennen. Zudem ermöglicht das eingesetzte Edge Computing eine lokale Datenverarbeitung nahezu in Echtzeit.~\acs{LPWAN} oder das vor allem genutzte \acs{LoRaWAN} bieten dem System außerdem eine hohe Robustheit im Gelände, da sie auch in Gebieten mit schlechter Netzabdeckung zuverlässig Daten übertragen können. Allerdings sind für eine effektive Abdeckung eines Gebietes hunderte Sensoren erforderlich. Dies führt zu einer erheblichen Steigerung der Kosten und somit zu einer starken Einschränkung der ökonomischen Effizienz bei einer flächendeckenden Anwendung~\cite{Chan2024,Kumar2023, Nodes.sh2024}.

\chapter{Identifikation technologischer Lücken}\label{ch:identifikationTechnologischerLuecken}

Basierend auf dem Vergleich anhand der fünf Kriterien wird deutlich, dass die technologischen Lücken nicht primär in der mangelnden Qualität der Systeme liegen, denn es existieren hochmoderne Satelliten, echtzeitfähige Funknetzwerke und hochentwickelte KI-Algorithmen. Vielmehr besteht die Herausforderung in der Kombination der Technologien, um die jeweiligen Stärken zu nutzen und die Schwächen zu kompensieren~\cite{Chan2024}.

Ein wesentlicher Aspekt stellt hier der Konflikt zwischen der Detektionslatenz und der Skalierbarkeit dar. Während \acs{IoT}-Sensoren eine extrem niedrige Latenz bieten, sind sie aufgrund der erforderlichen Anzahl an Sensorknoten und den damit verbundenen Infrastrukturkosten nicht flächendeckend einsetzbar. Satelliten hingegen bieten eine großflächige Abdeckung und skalieren global sowie effizient, können einen Brand aber erst dann wahrnehmen, wenn er groß genug ist, um das Blätterdach thermisch oder optisch zu durchdringen, was die Echtzeiterkennung oftmals limitiert. Optisch-terrestrische Systeme liegen in der Mitte, können aber durch die erforderliche Turm- und Netzinfrastruktur nicht flächendeckend eingesetzt werden~\cite{Allison2016}.

Ein weiterer Konflikt besteht zwischen der Fehlalarmresistenz und der visuellen Verifizierbarkeit. Ein \acs{IoT}-Sensor ist in der Lage, chemische Veränderungen ohne große Verzögerung zu erkennen, jedoch kann ohne eine visuelle Verifizierung nicht gewährleistet werden, dass es sich tatsächlich um einen Brand handelt~\cite{Kumar2023}. Außerdem fehlt durch das reine Vertrauen auf die Sensordaten das genaue Ausmaß der Flammen, die Ausbreitungsrichtung und die Kenntnis über betroffene Zufahrtswege. Im ungünstigsten Szenario führt dies zu einer Alarmierung auf Basis einer ungesicherten Verdachtslage, ohne dasss den Einsätzkräften eine fundierte taktische Lagebeurteilung zugänglich gemacht wird. Da nicht ersichtlich ist, wie groß der Waldbrand bereits ist, kann es passieren, dass zu viele oder zu wenige Einsatzkräfte alarmiert werden, was mit erhöhten Kosten oder verheerenden Verzögerungen bei den Löscharbeiten einhergeht. Satelliten und optisch-terrestrische Systeme können hingegen visuell verifizieren, sind aber durch Wolken oder topografische Hindernisse in ihrer Reaktionsgeschwindigkeit oftmals eingeschränkt~\cite{Bernhard2010}.

Außerdem besteht eine Lücke in der Robustheit der Systeme und der Kommunikationsinfrastruktur. Obwohl \acs{LoRaWAN}-Sensoren am Boden sehr robust und energieeffizient arbeiten, sind Gateways nötig, um die Daten via 4G, 5G oder Kabel an einen Server weiterzuleiten~\cite{Nodes.sh2024}. In sehr abgelegenen Gebieten fehlt diese notwendige Infrastruktur, wodurch die Alarmierung verzögert oder unmöglich wird. Satelliten und optisch-terrestrische Systeme sind hier weniger abhängig von der lokalen Infrastruktur, können aber durch atmosphärische Bedingungen oder topografische Hindernisse beeinträchtigt werden~\cite{VIIRSLexikon}.

\chapter{Ableitung eines multimodalen Konzepts}\label{ch:ableitungmultimodalenKonzep}

Um die identifizierten technologischen Lücken zu schließen und alle definierten Anforderungen vollumfänglich zu erfüllen, bedarf es eines multimodalen Frühwarnsystems, welches die Stärken der einzelnen Technologien kombiniert und die jeweiligen Schwächen kompensiert. Aktuelle Forschungsprojekte und kommerzielle Entwicklungen zeigen, wie terrestrische \acs{IoT}-Sensorik, autonome Flugsysteme und satellitengestützte Erdbeobachtung in einer hybriden Architektur zusammengeführt werden können.

Ein solches System könnte beispielsweise wie folgt aufgebaut sein:
Um den Konflikt zwischen der Robustheit der Systeme und der Kommunikationsinfrastruktur zu lösen, werden zunehmend terrestrische \acs{LPWAN}-Netzwerke mit Satellitentechnologie gekoppelt. Dryad Networks hat hierfür in Kooperation mit dem französischen Satellitenbetreiber Kinéis eine Technologie entwickelt, die die direkte Übertragung von Daten der \acs{LoRaWAN}-Sensoren an Satelliten realisiert. Diese \enquote{Direct-to-Satellite}-Verbindung sorgt dafür, dass Gateways Alarmmeldungen direkt an eine Satellitenkonstellation im Orbit senden, falls kein terrestrisches Mobilfunknetz zur Verfügung steht. Dies gewährleistet eine lückenlose und weltweite Alarmierungskette, selbst in unzugänglichen Gebirgsregionen oder entlang kritischer linearer Infrastrukturen wie Bahn- und Stromtrassen~\cite{DryadNetworks2024}.

Auch für die Lücke zwischen der Fehlalarmresistenz und der visuellen Verifizierbarkeit lässt sich durch eines multimodales Konzept eine Lösung finden. Die Sensoren am Boden sollen als automatische Auslöser (\enquote{Trigger}) agieren und unbemannte Luftfahrzeuge (\acs{UAV}s/Drohnen) zur visuellen Lagebeurteilung anfordern und autonom zum Detektionsort navigieren. Über ein latenzarmes 5G-Netz kann die Drohne so vor dem Eintreffen der ersten Einsatzkräfte hochauflösende optische und thermische Echtzeit-Videodaten an die Einsatzzentrale übermitteln. Die visuelle Verifikation erlaubt eine Einschätzung der Brandausbreitung und der betroffenen Zufahrtswege. Hierdurch wird eine gezielte Alarmierung der Einsatzkräfte gewährleistet und folglich können Kosten und Verzögerungen bei Löscharbeiten minimiert werden~\cite{Dauernheim2025, Laessig2025}.

Moderne Kleinsatelliten-Konstellationen bieten die Lösung für den Konflikt zwischen Detektionslatenz und Skalierbarkeit für Regionen, in denen eine flächendeckende \acs{IoT}-Sensorik unwirtschaftlich ist. Unternehmen wie OroraTech setzen hierbei auf Edge-Computing im Orbit und sogenannte \enquote{Thermal-Livestreams}~\cite{Seton2025}. 
Damit das multimodale Konzept auch praktisch umsetzbar ist, müssen alle Datenströme in einer zentralen Einsatzzentrale zusammengeführt werden. Im Forschungsprojekt \enquote{5G-Waldwächter} passiert das mittels Unterstützung einer KI über einen Zentralen Krisenmanagement-Server (\acs{ZKMS}). Dies ermöglicht auch Voraussagen zu erwarteten Brandausbreitungen, um Einsatzkräfte frühzeitig auf die zu erwartende Lage vorzubereiten~\cite{Laessig2025}. 





\part{Diskussion und Fazit}\label{pt:diskussionFazit}
\chapter{Diskussion}\label{ch:diskussion}

BLABLABLA


\chapter{Fazit und Ausblick}\label{ch:fazitAusblick}

BLABLABLA

\part{Appendix}\label{pt:appendix}
\printbibliography

\end{sloppypar}
\end{document}
